\pagestyle{plain}

\chapter{CAPÍTULO 2: Vulnerabilidade Biocultural em Saberes Agrícolas Tradicionais: uma Meta-análise Quantitativa}

\begin{flushright}
Catuxe Varjão de Santana Oliveira  \footnote[1]{Programa de Pós-graduação em Ciências da Propriedade Intelectual-PPGPI/Universidade Federal de Sergipe-UFS. *E-mail:  \href{mailto:catuxe@academico.ufs.br}{\color{blue}{\underline{catuxe@academico.ufs.br}}}}
\\[0.3cm]
Luiz Diego Vidal Santos \footnote[2]{Universidade Estadual de Feira de Santana - UEFS, Brasil. E-mail: \href{mailto:ldvsantos@uefs.br}{\color{blue}{\underline{ldvsantos@uefs.br}}}}
\\[0.3cm]
Paulo Roberto Gagliardi \footnotemark[1]
\end{flushright}

% resumo em português — bookmark suprimido para manter hierarquia do capítulo
{\renewcommand{\PRIVATEbookmarkthis}[1]{}%
\begin{resumoumacoluna}
A vulnerabilidade biocultural dos saberes e sistemas agrícolas tradicionais (SSAT) quilombolas permanece avaliada de forma fragmentada na literatura, impedindo a priorização de salvaguardas parametrizadas por evidência. Este estudo sintetizou, sob diretrizes PRISMA~2020, a resposta de seis dimensões funcionais de vulnerabilidade (erosão intergeracional $V_{1}$, complexidade biocultural $V_{2}$, singularidade territorial $V_{3}$, status de documentação $V_{4}$, vulnerabilidade jurídica $V_{5}$ e organização social $V_{6}$) a intervenções de salvaguarda em contextos quilombolas e análogos do Semiárido brasileiro, empregando razões de resposta log-transformadas ($lnRR$) em modelo hierárquico de efeitos mistos com heterogeneidade quantificada por $Q$ de Cochran e $I^{2}$, viés avaliado pelos testes de Egger e Begg e transferibilidade modelada por meta-regressão com congruência geográfica do Semiárido Nordeste~II. Os rankings de magnitude de efeito e o mapeamento de lacunas constituem arcabouço decisório para gestão de risco biocultural em territórios quilombolas e subsidiam a construção futura de um índice composto de salvaguarda.

\vspace{\onelineskip}
\noindent
\textbf{Palavras-chave}: Meta-análise quantitativa $\cdot$ Vulnerabilidade biocultural $\cdot$ Saberes agrícolas tradicionais quilombolas $\cdot$ Semiárido brasileiro.
\end{resumoumacoluna}}

\newpage

%% ============================================================
%% 1. INTRODUÇÃO
%% ============================================================
\section{Introdução}
\label{sec:intro_cap2}

A expansão de sistemas agrícolas intensivos sobre biomas de alta sociobiodiversidade sobrepõe-se a territórios de comunidades detentoras de saberes agroecológicos, gerando vetores diretos de erosão epistêmica nas fronteiras de uso do solo \cite{Altieri2004}. Reduções empíricas de 30 a 50\% no domínio de rotinas etnobotânicas entre gerações consecutivas confirmam a severidade dessa perda em comunidades submetidas à homogeneização produtiva \cite{ReyesGarcia2013}. 

Nesse mesmo processo de simplificação, o decaimento de variedades genéticas cultivadas compromete mecanismos de estabilização ambiental em agroecossistemas tradicionais \cite{Jarvis2008}, embora esses saberes constituam ativos intangíveis compatíveis com os critérios de vantagem competitiva sustentável formulados por \citeonline{Barney1991}. Apesar dessa relevância, tais ativos permanecem invisíveis perante instâncias regulatórias de proteção e monitoring, como IPHAN, INPI, CGen e FAO-GIAHS, perpetuando a vulnerabilidade territorial e epistêmica de comunidades quilombolas do semiárido brasileiro \cite{Santilli2009,Cunha2009,Goncalves2022}.

A vulnerabilidade biocultural desses sistemas não opera, contudo, como vetor unidimensional. Quando a proporção de jovens que dominam rotinas agroecológicas decresce entre coortes consecutivas, o que se perde não é apenas informação técnica, mas a infraestrutura cognitiva que sustenta a manutenção de variedades locais, a leitura de sinais edafoclimáticos e a adaptação de calendários de cultivo \cite{ReyesGarcia2013,GomezBaggethun2013}.

Essa erosão intergeracional desencadeia um efeito cascata sobre a complexidade biocultural dos agroecossistemas, pois, a simplificação do repertório de manejo reduz a riqueza de espécies cultivadas e elimina interações solo-planta cuja codificação depende de transmissão oral continuada \cite{Brush2004,Jarvis2008,Bastos2022}. 

O resultado é um sistema progressivamente homogeneizado, cuja perda de capacidade adaptativa compromete a resiliência frente a perturbações climáticas e econômicas, retroalimentando a própria erosão dos saberes que o sustentavam \cite{Tran2025,Aniah2019}.

Paralelamente, a exclusividade geográfica dos repertórios de manejo confere a cada território um grau de singularidade cuja perda é irreversível, uma vez que práticas endêmicas, desenvolvidas ao longo de gerações sob condições edafoclimáticas específicas, não são reconstituíveis por transferência horizontal entre comunidades \cite{Laborde2023,Ens2016}. 

A literatura reconhece essa singularidade, porém a avalia predominantemente por inventários qualitativos que não permitem comparação paramétrica entre territórios \cite{CalvetMir2016,LaRosa2021,Saiba2023}, impedindo a identificação de quais comunidades concentram os maiores riscos de perda irreversível.

A fragilidade se amplifica no plano institucional. A ausência de registros formais estruturados expõe os saberes a dois riscos convergentes, a invisibilidade perante políticas públicas de proteção e a vulnerabilidade a apropriação indevida por agentes externos, cenário agravado pela inexistência de instrumentos jurídicos efetivos que reconheçam o caráter coletivo e dinâmico desses ativos intangíveis \cite{Santilli2009}. Sem indicações geográficas, marcas coletivas ou registros de patrimônio imaterial que vinculem formalmente o saber ao território, a proteção depende exclusivamente da capacidade organizacional da própria comunidade, cuja vitalidade, expressa pela frequência de eventos de transmissão, pelo número de mestres de saberes ativos e pela existência de protocolos comunitários de governança, constitui a última linha de defesa contra a desagregação do sistema \cite{FalsBorda1991,Ostrom1990}. A degradação dessa organização social fecha o ciclo de retroalimentação, acelerando a erosão intergeracional que originou todo o processo \cite{Agrawal2002}.

Diante da fragmentação das abordagens disponíveis na literatura, a meta-análise quantitativa guiada por diretrizes PRISMA~2020 constitui a ferramenta metodológica adequada para sintetizar tamanhos de efeito padronizados de múltiplos estudos empíricos, modelar heterogeneidade entre contextos e identificar moderadores da variabilidade nos resultados \cite{Page2021,Borenstein2009,Pustejovsky2022}.

A hipótese direcional postula que dimensões vinculadas à preservação in situ e à transmissão intergeracional apresentarão magnitudes de efeito superiores às dimensões de salvaguarda normativa ou documental, refletindo assimetrias na inércia ecológica de cada componente biocultural. A meta-análise integra uma arquitetura mais ampla de salvaguarda de SSAT quilombolas, orientada a converter conhecimento tácito em parâmetros mensuráveis e gerar documentação auditável para processos de salvaguarda patrimonial e registro de propriedade intelectual coletiva.

Nesse contexto, o objetivo deste estudo é determinar quais dimensões de vulnerabilidade biocultural exercem maior magnitude de efeito e maior sensibilidade a intervenções de salvaguarda em SSAT quilombolas e contextos análogos do Semiárido brasileiro, hierarquizando as dimensões segundo magnitude de efeito e heterogeneidade e mapeando sistematicamente as lacunas de evidência que orientarão investigações complementares.

%% ============================================================
%% 2. MATERIAIS E MÉTODOS
%% ============================================================

\section{Materiais e Métodos}
\label{sec:methods_cap2}

\subsection{Registro e conformidade PRISMA 2020}
\label{sec:prisma_cap2}

Este estudo foi conduzido em conformidade com as diretrizes PRISMA~2020 para revisões sistemáticas e meta-análises \cite{Page2021}. 

A Figura~\ref{fig:prisma-flow_cap2} apresenta o diagrama de fluxo PRISMA~2020, que sintetiza o processo de identificação, triagem, deduplicação e seleção dos 48 estudos elegíveis a partir das duas bases consultadas.

\begin{figure}[!ht]
\centering
\includegraphics[width=0.88\textwidth]{CAPITULOS/ARTIGO_2/figuras/prisma_flowdiagram_artigo2_en.pdf}
\caption{Diagrama de fluxo PRISMA 2020 mapeando a busca sistemática, restrições de triagem e a desduplicação entre bancos de dados.}
\label{fig:prisma-flow_cap2}
\end{figure}

O protocolo foi registrado prospectivamente na plataforma PROSPERO (número de registro CRD42024xxxxxx; \url{https://www.crd.york.ac.uk/prospero/}). A condução seguiu adicionalmente as recomendações do Cochrane Handbook for Systematic Reviews of Interventions \cite{Higgins2019} para extração de dados, avaliação de risco de viés e síntese quantitativa.

\subsection{Framework conceitual}
\label{sec:framework_cap2}

A vulnerabilidade biocultural de Saberes e Sistemas Agrícolas Tradicionais (SSAT) foi modelada como construto latente hexadimensional, fundamentado na convergência entre a taxonomia de riscos de erosão de conhecimento tradicional \cite{Agrawal2002}, os indicadores de resiliência de sistemas socioecológicos \cite{Folke2016} e dimensões de vulnerabilidade identificadas na literatura. O framework postula que a vulnerabilidade biocultural agregada resulta da interação de seis dimensões funcionais ($V_1$--$V_6$), cada uma quantificável por indicadores proxy extraíveis da literatura empírica, e que intervenções de salvaguarda produzem efeitos mensuráveis e diferenciados sobre essas dimensões. A magnitude relativa desses efeitos fornece a hierarquização empírica das dimensões de vulnerabilidade, cujos rankings e parâmetros de heterogeneidade poderão subsidiar a construção de instrumentos compostos de avaliação em estudos subsequentes. O mapeamento sistemático da densidade de evidência por dimensão orienta a identificação de lacunas e a elicitação de dados primários em investigações complementares. A ausência de síntese quantitativa capaz de hierarquizar a magnitude relativa das forças que atuam sobre cada dimensão e sua sensibilidade a intervenções de salvaguarda compromete a capacidade de priorização parametrizada por evidência, lacuna que justifica a arquitetura analítica detalhada nas subseções seguintes.

\subsection{Critérios de elegibilidade (PICOS)}
\label{sec:picos_cap2}

Os critérios de elegibilidade foram definidos pelo framework PICOS (Tabela~\ref{tab:picos_cap2}) para assegurar transparência, replicabilidade e comparabilidade dos estudos incluídos.

\begin{table}[htbp]
\caption{Framework PICOS (critérios de elegibilidade para a meta-análise de dimensões de vulnerabilidade biocultural).}
\label{tab:picos_cap2}
\centering
\small
\begin{tabular}{p{2.5cm}p{10.5cm}}
\toprule
\textbf{Elemento} & \textbf{Descrição} \\
\midrule
P (População) & Saberes e Sistemas Agrícolas Tradicionais (SSAT) mantidos por comunidades quilombolas do Semiárido brasileiro (bioma Caatinga) como população primária, e comunidades indígenas, camponesas ou ribeirinhas de outras regiões tropicais e subtropicais como amostra comparativa de controle. \\
I (Intervenção) & Práticas de salvaguarda, documentação estruturada, registro patrimonial, proteção jurídica, governança comunitária de saberes, ou programas de revitalização de transmissão intergeracional. \\
C (Comparação) & Comunidades ou sistemas sem intervenções formais de salvaguarda, com documentação convencional não padronizada, ou condições pré-intervenção (delineamentos antes-depois). \\
O (Desfechos) & Indicadores quantitativos de vulnerabilidade biocultural operacionalizados como proxy das seis dimensões ($V_1$--$V_6$), incluindo proporção de jovens praticantes, índices de diversidade de espécies, indicadores de endemismo de práticas, taxas de registro formal, métricas de proteção jurídica e indicadores de vitalidade organizacional (Tabela~\ref{tab:proxies_cap2}). \\
S (Delineamento) & Estudos empíricos comparativos (experimentais, quasi-experimentais, transversais comparativos com grupo de comparação identificável, ou longitudinais antes-depois) publicados entre 2016 e 2026, em inglês, português ou espanhol, reportando dados quantitativos (médias, tamanhos amostrais e medidas de variância) ou evidência semi-quantitativa convertível (p-values com tamanhos amostrais, estatísticas de teste ANOVA/Kruskal-Wallis, ou julgamentos direcionais codificáveis em escala ordinal), conforme o protocolo de integração multi-tier descrito adiante. \\
\bottomrule
\end{tabular}
\end{table}

A operacionalização do componente Outcomes admitiu cinco tiers de evidência em ordem decrescente de precisão, desde indicadores contínuos com variâncias completas (Tier~1) até julgamentos qualitativos direcionais convertidos em $lnRR$ sintéticos (Tier~4). A janela temporal (2016--2026) concentra o maior volume de produção sobre vulnerabilidade biocultural e salvaguarda de conhecimentos tradicionais. Foram incluídos delineamentos transversais comparativos (comunidades com vs.\ sem intervenção) e longitudinais antes-depois, reconhecendo que intervenções de salvaguarda biocultural raramente permitem randomização controlada \cite{Shadish2002}.

\subsection{Estratégia de busca e bases de dados}
\label{sec:search_cap2}

Registros bibliográficos foram sistematicamente recuperados de duas bases de dados indexadas, Scopus e Web of Science Core Collection. A estratégia de busca combinou descritores em inglês, português e espanhol organizados em três blocos semânticos conectados por operador AND, com variações internas conectadas por OR.

O bloco populacional incluiu os termos \textit{traditional agricultural system}, \textit{traditional ecological knowledge}, \textit{indigenous farming}, \textit{quilombola}, \textit{smallholder}, \textit{agrobiodiversity}, \textit{agroecology} e equivalentes em português e espanhol. O bloco de intervenção e exposição compreendeu \textit{safeguard*}, \textit{documentation}, \textit{heritage registration}, \textit{legal protection}, \textit{intangible heritage}, \textit{biocultural conservation}, \textit{community governance} e equivalentes. O bloco de desfechos incluiu \textit{vulnerability}, \textit{erosion}, \textit{knowledge loss}, \textit{intergenerational transmission}, \textit{biodiversity index}, \textit{cultural complexity}, \textit{documentation status} e equivalentes.

As buscas foram complementadas por rastreamento de referências dos estudos incluídos (\textit{backward snowballing}) e citações dos estudos incluídos (\textit{forward snowballing} via Google Scholar), para minimizar omissão de estudos relevantes não indexados nas bases primárias. Todos os registros foram exportados em formato BibTeX, unificados em base de dados consolidada e submetidos a remoção automática de duplicatas mediante o pacote \textit{revtools} do R \cite{Westgate2019}.

\subsection{Triagem e seleção}
\label{sec:screening_cap2}

Previamente à triagem formal, os dois revisores conduziram exercício piloto de calibração com amostra aleatória de 50 registros para padronizar a interpretação dos critérios de elegibilidade e reduzir discordâncias sistemáticas. Discrepâncias identificadas no piloto foram discutidas até consenso operacional, e a concordância resultante superou o limiar de concordância substancial ($\kappa \geq 0{,}61$), habilitando o início da triagem em escala.

A triagem formal foi conduzida em duas fases por dois revisores independentes. Na fase de títulos e resumos, todos os registros foram avaliados contra os critérios PICOS, com discordâncias resolvidas por consenso ou consulta a terceiro revisor. A concordância interavaliadores na triagem formal foi quantificada pelo coeficiente kappa de Cohen ($\kappa$), com limiar de concordância substancial ($\kappa \geq 0{,}61$) conforme critérios de \citeonline{Landis1977}. Na fase de texto completo, os registros pré-selecionados foram avaliados integralmente para verificação de adequação aos critérios de elegibilidade, disponibilidade de dados quantitativos para extração e ausência de sobreposição amostral com outros estudos incluídos.

Revisões narrativas, capítulos de livro, editoriais e relatos de experiência sem qualquer dado comparativo foram excluídos. Publicações com evidência exclusivamente qualitativa, porém contendo julgamentos direcionais sobre mudança em indicadores de vulnerabilidade biocultural (aumento, estabilidade ou redução), foram retidas e classificadas como Tier~3 ou Tier~4 conforme o protocolo de integração multi-tier, assegurando representatividade do espectro metodológico da literatura interdisciplinar. Para estudos que reportaram dados em formato insuficiente (médias sem variância, ou dados exclusivamente gráficos), foi tentado contato com os autores correspondentes antes da reclassificação ao tier correspondente. O fluxo de seleção seguiu o diagrama PRISMA~2020 \cite{Page2021}, documentando registros em cada etapa, razões de exclusão e o tier de evidência atribuído.

\subsection{Classificação funcional das dimensões de vulnerabilidade ($V_1$--$V_6$) e variáveis proxy}
\label{sec:dimensions_cap2}

Os indicadores extraídos dos estudos primários foram organizados em seis dimensões funcionais de vulnerabilidade biocultural, derivadas da convergência entre a revisão de escopo e a taxonomia de riscos de erosão do conhecimento tradicional \cite{Agrawal2002}. Para cada dimensão, foram definidos indicadores proxy operacionalizáveis, unidades de medida esperadas e a disponibilidade estimada de estudos empíricos quantitativos na literatura (Tabela~\ref{tab:proxies_cap2}).

\begin{table}[htbp]
\caption{Dimensões de vulnerabilidade biocultural ($V_1$--$V_6$), variáveis proxy operacionalizáveis, unidades de medida e disponibilidade estimada de evidência quantitativa para cálculo de $lnRR$.}
\label{tab:proxies_cap2}
\centering
\footnotesize
\begin{tabular}{p{0.8cm}p{2.2cm}p{4.5cm}p{2.5cm}p{1.5cm}}
\toprule
\textbf{Dim.} & \textbf{Denominação} & \textbf{Variáveis proxy (indicadores extraíveis)} & \textbf{Unidade de medida} & \textbf{Disp.} \\
\midrule
$V_1$ & Erosão Intergeracional & Proporção de jovens ($<$ 35 anos) com domínio do saber; taxa de transmissão ativa (eventos/ano); índice de conhecimento etnobotânico por faixa etária & \%, eventos$\cdot$ano$^{-1}$, escore & Alta \\
$V_2$ & Complexidade Biocultural & Riqueza de espécies manejadas (Shannon $H'$, Simpson $1-D$); número de variedades locais mantidas; número de interações solo-planta codificadas no saber & índice, contagem & Alta \\
$V_3$ & Singularidade Territorial & Índice de beta-diversidade de práticas entre comunidades (Jaccard, Bray-Curtis); índice de endemismo de variedades locais; proporção de práticas exclusivas (\%) & índice, \% & Baixa \\
$V_4$ & Status de Documentação & Número de saberes com registro formal (inventários, audiovisual); completude de fichas técnicas (\%); taxa de digitalização de acervos (\%) & contagem, \% & Moderada \\
$V_5$ & Vulnerabilidade Jurídica & Número de mecanismos formais de proteção existentes (IG, marca coletiva, registro IPHAN); existência de protocolo comunitário; lacuna regulatória & contagem, binário, ordinal & Baixa \\
$V_6$ & Organização Social & Número de mestres de saberes ativos; frequência de reuniões comunitárias de transmissão (eventos/mês); vitalidade de protocolos comunitários de governança & contagem, eventos$\cdot$mês$^{-1}$, escore & Moderada \\
\bottomrule
\end{tabular}
\end{table}

As dimensões $V_1$ (Erosão Intergeracional) e $V_2$ (Complexidade Biocultural) operacionalizam, respectivamente, o risco de descontinuidade na transmissão vertical do conhecimento e a riqueza de interações ecológicas codificadas nos sistemas tradicionais. $V_1$ captura proporções de jovens com domínio do saber e índices de conhecimento etnobotânico estratificados por faixa etária, com alta disponibilidade de estudos comparativos em etnobotânica \cite{ReyesGarcia2013,GomezBaggethun2013}. $V_2$ emprega índices de diversidade (Shannon $H'$, Simpson $1-D$) e contagens de variedades locais, sustentada pela literatura de agrobiodiversidade e conservação \textit{in situ} \cite{Jarvis2008,Brush2004}. A dimensão $V_3$ (Singularidade Territorial) captura a exclusividade geográfica das práticas via índices de beta-diversidade (Jaccard, Bray-Curtis) e taxa de endemismo de variedades locais, admitindo como proxies complementares os índices de dissimilaridade calculáveis a partir de inventários de $V_1$ e $V_2$.

As dimensões institucionais $V_4$ (Status de Documentação) e $V_5$ (Vulnerabilidade Jurídica) mensuram, respectivamente, o grau de registro formal dos saberes e a exposição a apropriação indevida por ausência de instrumentos de proteção (indicações geográficas, marcas coletivas, registros de patrimônio imaterial). Estudos avaliativos de programas patrimoniais (IPHAN, UNESCO, FAO-GIAHS) fornecem dados de completude pré/pós-intervenção, e relatórios institucionais foram considerados elegíveis quando contiveram grupo de comparação. A dimensão $V_6$ (Organização Social para Transmissão) quantifica a vitalidade das estruturas comunitárias de governança do conhecimento (mestres de saberes ativos, frequência de eventos de transmissão, protocolos comunitários), com dados oriundos de surveys de capital social e ação coletiva \cite{Ostrom1990,Pretty2003}.

\subsection{Diagnóstico de disponibilidade de evidência e estratégia analítica por dimensão}
\label{sec:diagnosis_cap2}

A estratégia analítica foi diferenciada por dimensão, adequando a complexidade do modelo à densidade de evidência disponível. A Tabela~\ref{tab:viability_cap2} sintetiza a estratégia adotada para cada dimensão.

\begin{table}[htbp]
\caption{Diagnóstico de disponibilidade de evidência e estratégia analítica por dimensão de vulnerabilidade.}
\label{tab:viability_cap2}
\centering
\small
\begin{tabular}{p{0.8cm}p{2.2cm}p{1.2cm}p{2.0cm}p{5.5cm}}
\toprule
\textbf{Dim.} & \textbf{Denominação} & \textbf{$k$ esp.} & \textbf{Status} & \textbf{Estratégia analítica} \\
\midrule
$V_1$ & Erosão Interg. & $\geq$ 15 & Confirmatória & Meta-análise completa (forest plot, meta-regressão, subgrupos, viés de publicação) \\
$V_2$ & Complexidade & $\geq$ 15 & Confirmatória & Meta-análise completa \\
$V_3$ & Singularidade & 5--10 & Exploratória & Efeito agregado com IC amplo; sem análise de subgrupos; sensibilidade \textit{leave-one-out}; reportada como exploratória \\
$V_4$ & Documentação & 8--12 & Confirmatória & Meta-análise completa; subgrupos condicionais ($k \geq 10$) \\
$V_5$ & Vuln. Jurídica & $<$ 8 & Exploratória & Efeito agregado se $k \geq 3$; síntese narrativa se $k < 3$; reportada como exploratória \\
$V_6$ & Org. Social & 10--15 & Confirmatória & Meta-análise completa; subgrupos condicionais ($k \geq 10$) \\
\bottomrule
\end{tabular}
\end{table}

Conforme \citeonline{Borenstein2009}, a complexidade do modelo foi proporcionada à evidência disponível. Dimensões com $k \geq 10$ receberam meta-análise completa (subgrupos e meta-regressão multivariada), enquanto $V_3$ e $V_5$ ($k < 10$) foram tratadas com estratégia exploratória (efeito agregado e sensibilidade \textit{leave-one-out}), cujas estimativas são reportadas com cautela. O critério mínimo para inclusão na síntese quantitativa foi fixado em $k \geq 3$ estudos independentes.

\subsection{Protocolo de extração de dados}
\label{sec:extraction_cap2}

A extração de dados seguiu protocolo estruturado em três blocos, com formulário eletrônico padronizado preenchido por dois revisores independentes e discordâncias resolvidas por consenso.

O bloco de metadados do estudo registrou autor, ano de publicação, DOI, país, bioma ou região biogeográfica (semiárido, cerrado, amazônia, mata atlântica, savana, floresta tropical, outro), tipo de comunidade (quilombola, indígena, camponesa, ribeirinha, outra), tipo de intervenção (documentação audiovisual, registro patrimonial formal, proteção jurídica ativa, governança comunitária, programa de revitalização de transmissão, nenhuma/controle), duração da intervenção em anos e delineamento (experimental, quasi-experimental, transversal comparativo, longitudinal antes-depois).

O bloco de dados para cálculo de tamanho de efeito registrou, para cada comparação pareada, a dimensão funcional atribuída ($V_1$--$V_6$), o indicador proxy específico (conforme Tabela~\ref{tab:proxies_cap2}), a unidade de medida, a média sob intervenção ($\bar{X}_T$), o desvio-padrão sob intervenção ($SD_T$), o tamanho amostral sob intervenção ($n_T$), a média sob controle ($\bar{X}_C$), o desvio-padrão sob controle ($SD_C$) e o tamanho amostral controle ($n_C$). Quando os estudos reportaram erro-padrão ($EP$), a conversão para desvio-padrão foi realizada por
%
\begin{equation}
SD = EP \times \sqrt{n}
\label{eq:sd-ep_cap2}
\end{equation}

\noindent Quando apenas intervalos de confiança foram disponíveis, médias e desvios-padrão foram estimados conforme \citeonline{Betancur2022}. Dados apresentados exclusivamente em figuras foram digitalizados mediante WebPlotDigitizer v4.3 \cite{Rohatgi2020}. Quando estudos reportaram resultados para múltiplas comunidades, práticas ou períodos, os dados foram extraídos separadamente para cada condição, preservando a resolução e evitando que efeitos de tratamento fossem obscurecidos por agregação.

O bloco de moderadores registrou as variáveis moderadoras pré-definidas para meta-regressões. O moderador \textit{tipo de intervenção} diferenciou quatro categorias (documentação, registro patrimonial, proteção jurídica, governança comunitária). O moderador \textit{região biogeográfica} controlou variabilidade edafoclimática e cultural (semiárido, cerrado, amazônia, mata atlântica, outro tropical, outro subtropical). Um moderador derivado, \textit{congruência com Semiárido NE~II}, foi codificado como variável binária (1 = estudo conduzido em comunidades quilombolas ou no bioma Caatinga do Território de Identidade Semiárido Nordeste II; 0 = demais contextos), permitindo quantificar a transferibilidade dos tamanhos de efeito para o contexto-alvo da arquitetura de salvaguarda.

O moderador \textit{tempo desde intervenção} foi registrado como variável contínua (anos) e categorizado em três níveis (curto prazo $<$ 3~anos, médio prazo 3--10~anos, longo prazo $>$ 10~anos) para análises de subgrupo. O moderador \textit{tipo de comunidade} diferenciou quilombolas, indígenas, camponesas e ribeirinhas. O escore de qualidade metodológica (NOS) foi registrado como moderador contínuo para verificar associação entre qualidade dos estudos e magnitude dos efeitos.

\subsection{Protocolo de integração de evidência multi-tier}
\label{sec:multitier_cap2}

Os 48 estudos elegíveis foram classificados em cinco tiers de evidência segundo a completude dos dados reportados, seguindo protocolo de integração multi-tier com conversões calibradas e propagação explícita de incerteza. Estudos com médias, desvios-padrão e tamanhos amostrais para ambos os grupos (Tier~1) permitiram cálculo direto de $lnRR$ e $v_i$ conforme Equações~\ref{eq:lnrr_cap2} e~\ref{eq:vi_cap2}.

Quando apenas p-values numéricos e tamanhos amostrais eram identificáveis (Tier~2a), o tamanho de efeito foi obtido pela cadeia $p \rightarrow z \rightarrow d \rightarrow lnOR \rightarrow lnRR$, operacionalizada pela sequência de conversões expressa nas Equações~\ref{eq:z-conv_cap2}--\ref{eq:lnrr-conv_cap2}.
%
\begin{align}
z &= \Phi^{-1}(1 - p/2) \label{eq:z-conv_cap2} \\
d &= z \cdot \sqrt{\frac{1}{n_T} + \frac{1}{n_C}} \label{eq:d-conv_cap2} \\
lnOR &= \frac{d \cdot \pi}{\sqrt{3}} \label{eq:lnor-conv_cap2} \\
lnRR &= lnOR - \ln\!\bigl(1 - p_0 + p_0 \cdot e^{lnOR}\bigr) \label{eq:lnrr-conv_cap2}
\end{align}

\noindent A Equação~\ref{eq:z-conv_cap2} converte o p-value bilateral no quantil $z$ da distribuição normal padrão, e a Equação~\ref{eq:d-conv_cap2} traduz $z$ na diferença de médias padronizada $d$ \cite{Borenstein2009}. A Equação~\ref{eq:lnor-conv_cap2} transforma $d$ em log-odds ratio pela conversão de \citeonline{Hasselblad1995}, e a Equação~\ref{eq:lnrr-conv_cap2} converte o $lnOR$ resultante em $lnRR$ pela transformação de \citeonline{Zhang1998} com prevalência basal $p_0 = 0{,}5$ como prior não-informativo. Quando o coeficiente de variação ($CV$) mediano dos estudos Tier~1 na mesma dimensão estava disponível, a conversão alternativa
%
\begin{equation}
lnRR \approx d \cdot CV_{ref}
\label{eq:lnrr-cv_cap2}
\end{equation}

\noindent proposta por \citeonline{Lajeunesse2011} foi empregada por sua maior precisão.

A mesma cadeia foi aplicada a estudos que reportaram estatísticas de teste sem p-values explícitos (Tier~2b), adotando thresholds conservadores ($p = 0{,}05$ para efeitos moderados/fortes, $p = 0{,}10$ para efeitos fracos). Estudos que apresentaram comparações numéricas tabuladas sem estatísticas de teste (Tier~3) e estudos com conclusões exclusivamente narrativas (Tier~4) foram submetidos à codificação ordinal independente por dois revisores, com inflação adicional de variância ($\times 1{,}5$) no caso qualitativo.

A incerteza associada a cada conversão foi propagada inflando o erro-padrão conforme
%
\begin{equation}
SE_{adj} = SE \cdot \sqrt{1 + \sigma^2_{conv}}
\label{eq:se-adj_cap2}
\end{equation}

\noindent onde $\sigma_{conv}$ assume valores crescentes por tier ($0{,}15$ para Tier~2a, $0{,}25$ para Tier~2b, $0{,}35$ para Tier~3 e $0{,}50$ para Tier~4), assegurando ponderação proporcional à precisão epistemológica de cada registro no modelo hierárquico. A robustez dos resultados à inclusão de registros convertidos foi avaliada por meta-regressão com a variável indicadora \textit{evidência direta} (Tier~1 vs.~demais) como moderador, e por análise de sensibilidade restrita ao subconjunto Tier~1.

\subsection{Codificação ordinal e conversão para $lnRR$ sintético}
\label{sec:coding_cap2}

Para os registros classificados como Tier~3 e Tier~4, dois revisores independentes codificaram cada comparação em duas dimensões, seguindo protocolo de codificação com calibração prévia ($\kappa \geq 0{,}61$). A dimensão de direção do efeito atribuiu valores $+1$ (vulnerabilidade agravada), $0$ (neutro) ou $-1$ (vulnerabilidade reduzida) com base na interpretação do texto original. A dimensão de intensidade classificou a magnitude percebida em fraca (1), moderada (2) ou forte (3) com base na ênfase narrativa, significância estatística reportada (quando disponível) e amplitude da mudança descrita.

Os escores ordinais foram convertidos em $lnOR$ sintéticos pelo mapeamento a quantis fixos da distribuição logit-normal com variância $\pi^2/3 \approx 3{,}29$ \cite{Hasselblad1995}, resultando em $lnOR$ de $-0{,}405$ (fraco), $+0{,}405$ (moderado) e $+1{,}386$ (forte), multiplicados pela direção codificada. A conversão subsequente de $lnOR$ para $lnRR$ empregou a transformação de \citeonline{Zhang1998} com $p_0 = 0{,}5$, e a variância do $lnRR$ sintético incorporou tanto a variância fixa logística ($\pi^2/3$) quanto a inflação por incerteza de conversão.

Como verificação de robustez, implementou-se um modelo bayesiano hierárquico em \textit{brms} \cite{Burkner2017} onde os $lnRR$ dos Tiers~1--2 entraram como observações gaussianas com erro-padrão conhecido, enquanto os julgamentos ordinais dos Tiers~3--4 entraram via família cumulativa (logit), compartilhando o mesmo parâmetro $\theta_{dim}$ no nível superior do modelo. Essa formulação permite que evidência qualitativa funcione como prior empírico atualizável por dados quantitativos, preservando coerência interna sem inflação de graus de liberdade e convergindo para o estimador frequencista quando $k$ é suficiente \cite{Rover2020}.

\subsection{Avaliação de qualidade metodológica}
\label{sec:nos_cap2}

A qualidade metodológica dos estudos incluídos foi avaliada mediante versão adaptada da Newcastle-Ottawa Scale (NOS) para estudos observacionais \cite{Wells2000}, ajustada ao contexto de avaliação de intervenções de salvaguarda biocultural. A escala adaptada compreende três domínios com pontuação máxima de 9 estrelas.

O domínio Seleção (máximo 4 estrelas) avaliou a representatividade da amostra comunitária, a seleção do grupo de comparação, a definição operacional da intervenção de salvaguarda e a demonstração de que o desfecho não estava presente no início do estudo (para delineamentos longitudinais). O domínio Comparabilidade (máximo 2 estrelas) avaliou se os grupos foram comparáveis em relação a fatores de confusão (tamanho da comunidade, acesso a infraestrutura, escolaridade média, tempo de reconhecimento formal) e se controles estatísticos foram aplicados. O domínio Desfechos (máximo 3 estrelas) avaliou se os indicadores foram mensurados com instrumentos validados ou padronizados, se o seguimento foi suficiente para capturar efeitos e se a taxa de atrito foi reportada e aceitável ($<$ 20\%).

Cada estudo recebeu pontuação total (0--9), classificada em alta qualidade (7--9 estrelas), moderada (4--6) e baixa (0--3). A pontuação NOS foi utilizada como moderador contínuo na meta-regressão para verificar associação entre qualidade e magnitude de efeito, e como critério de análise de sensibilidade (restrição a estudos de alta qualidade).

\subsection{Tratamento de dados faltantes e estrutura hierárquica}
\label{sec:missing_cap2}

Estudos que reportaram médias sem medidas de variância ($SD$ ou $EP$) foram retidos para evitar viés de exclusão, conforme recomendado por \citeonline{Burgess2013}. Valores de variância faltantes foram estimados por imputação múltipla via equações encadeadas (MICE), implementada no pacote \textit{mice} do R \cite{vanBuuren2011}, utilizando Árvores de Classificação e Regressão (CART) como método univariado para gerar 20 conjuntos imputados. As estimativas finais foram combinadas segundo as regras de Rubin, que ajustam as variâncias pela incerteza da imputação \cite{Rubin1987}. Análises de sensibilidade compararam os resultados com e sem imputação para verificar robustez. 

A base de dados foi estruturada hierarquicamente em três níveis para controlar dependências estatísticas decorrentes de múltiplas comparações extraídas de um mesmo estudo \cite{Pustejovsky2022}. O Nível~1 (observação) representou cada comparação pareada entre intervenção e controle, vinculada a um indicador proxy específico, com $lnRR$ e variância amostral calculados. O Nível~2 (estudo) agrupou observações provenientes da mesma publicação, com efeitos aleatórios capturando condições compartilhadas (contexto geográfico, equipe de pesquisa, protocolo de coleta). O Nível~3 (meta-análise) incorporou efeitos fixos (moderadores) e a variância residual entre estudos ($\tau^2$), estimados por REML. 

Essa estrutura evita a inflação artificial da precisão que resultaria de tratar observações dependentes como independentes \cite{Cheung2014}. A matriz de variância-covariância dos efeitos intra-estudo ($\mathbf{V}_j$) foi construída assumindo correlação constante $\rho$ entre pares de tamanhos de efeito provenientes do mesmo estudo, conforme o modelo de efeitos correlados de \citeonline{Pustejovsky2022}, com valor de referência $\rho = 0{,}8$. Análises de sensibilidade variaram $\rho$ no intervalo $[0{,}5;\; 1{,}0]$ para avaliar a robustez das estimativas globais e dos intervalos de confiança à especificação da correlação intra-estudo. A estimação robusta de variância (\textit{robust variance estimation}, RVE) com correção de pequena amostra (CR2) foi aplicada via pacote \textit{clubSandwich} (versão $\geq$ 0.5.5) para produzir erros-padrão válidos independentemente da especificação exata de $\rho$ \cite{Pustejovsky2022}.

A heterogeneidade global foi quantificada pelo teste $Q$ de Cochran e pelo índice $I^2$ (Equação~\ref{eq:i2_cap2}), interpretado segundo os limiares de \citeonline{Higgins2003}, em que $I^2 < 25\%$ indica baixa heterogeneidade, 25--75\% moderada e $> 75\%$ alta.

\begin{equation}
I^2 = \frac{Q - (k-1)}{Q} \times 100
\label{eq:i2_cap2}
\end{equation}

\subsection{Análise estatística}
\label{sec:stats_cap2}

O logaritmo da razão de resposta ($lnRR$) foi adotado como métrica de tamanho de efeito por sua interpretabilidade proporcional e menor sensibilidade a \textit{outliers} comparado à diferença de médias padronizada \cite{Hedges1999}, sendo calculado conforme Equação~\ref{eq:lnrr_cap2}.

\begin{equation}
lnRR = \ln\left(\frac{\bar{X}_{T}}{\bar{X}_{C}}\right)
\label{eq:lnrr_cap2}
\end{equation}

\noindent Valores de $lnRR > 0$ indicam melhoria do indicador sob intervenção de salvaguarda (aumento de diversidade, transmissão, documentação), enquanto $lnRR < 0$ indicam deterioração. Para facilitar a interpretação, os $lnRR$ foram retrotransformados em variação percentual pela expressão $(e^{lnRR} - 1) \times 100$. A variância amostral associada a cada $lnRR$ foi estimada conforme Equação~\ref{eq:vi_cap2} \cite{Hedges1999}.

\begin{equation}
v_i = \frac{SD_T^2}{n_T \cdot \bar{X}_T^2} + \frac{SD_C^2}{n_C \cdot \bar{X}_C^2}
\label{eq:vi_cap2}
\end{equation}

O modelo hierárquico de efeitos mistos foi ajustado com estimador de máxima verossimilhança restrita (REML) para estimar a variância entre estudos ($\tau^2$), conforme Equação~\ref{eq:model_cap2}.

\begin{equation}
lnRR_{ij} = \mu + \boldsymbol{\beta} \cdot \mathbf{X}_{ij} + u_j + \epsilon_{ij}
\label{eq:model_cap2}
\end{equation}

\noindent onde $lnRR_{ij}$ é o tamanho de efeito da observação $i$ no estudo $j$, $\mu$ é o intercepto (efeito médio global ou por dimensão), $\boldsymbol{\beta} \cdot \mathbf{X}_{ij}$ é o vetor de moderadores fixos, $u_j \sim N(0, \tau^2)$ é o efeito aleatório de estudo e $\epsilon_{ij}$ é o erro residual.

O intervalo de predição complementou o IC~95\% para indicar a faixa esperada de efeitos em estudos futuros \cite{IntHout2016}. Para dimensões com $k < 15$, aplicou-se adicionalmente a correção de Hartung-Knapp-Sidik-Jonkman (HKSJ), que emprega a distribuição $t$ com $k-1$ graus de liberdade, produzindo cobertura mais conservadora \cite{Higgins2019}. Estimativas de $\tau^2$ foram comparadas entre os métodos DerSimonian-Laird e REML para avaliar robustez à escolha do estimador \cite{Borenstein2009}.

A identificação de estudos atípicos ou influentes foi conduzida por resíduos estudentizados com correção de Bonferroni ($\alpha = 0{,}05$) e pela distância de Cook, com limiar $D_i > 4/k$ \cite{Viechtbauer2010}. Análises de sensibilidade \textit{leave-one-out} foram conduzidas para todas as dimensões, e a técnica \textit{leave-one-study-out} (excluindo todas as observações de um estudo de cada vez) foi aplicada para avaliar a dependência dos resultados a publicações individuais.

O viés de publicação foi avaliado por inspeção visual de funnel plots (dispersão de $lnRR$ vs.\ erro-padrão), complementada pelo teste de Begg e Mazumdar \cite{Begg1994} e pelo teste de Egger \cite{Sterne2005}. Para dimensões com $k < 10$, os testes foram reportados mas interpretados com cautela devido ao baixo poder estatístico, conforme \citeonline{Sterne2011}. O método \textit{trim-and-fill} de Duval e Tweedie foi aplicado como análise suplementar para estimar estudos potencialmente faltantes e o efeito ajustado \cite{Duval2000}.

Meta-regressões multivariadas com moderadores fixos foram conduzidas para cada dimensão com $k \geq 10$, testando se tipo de intervenção, região biogeográfica, tempo desde intervenção, tipo de comunidade e qualidade metodológica (NOS) explicaram proporção significativa da heterogeneidade residual. O teste omnibus ($Q_M$) avaliou a significância conjunta dos moderadores, e coeficientes individuais foram interpretados com IC~95\%. Para dimensões com $k < 10$, moderadores foram avaliados univariadamente como análise exploratória, sem ajuste para multiplicidade. Análises de subgrupo foram conduzidas para moderadores categóricos com $k \geq 10$, utilizando o teste $Q_{between}$ e reportando heterogeneidade residual ($I^2_{within}$) dentro de cada subgrupo.

Todas as análises foram conduzidas em R (versão 4.3+) com os pacotes \textit{metafor} \cite{Viechtbauer2010} (modelos hierárquicos, meta-regressão), \textit{clubSandwich} \cite{Pustejovsky2022} (RVE), \textit{meta} (forest plots, subgrupos), \textit{dmetar} (influência, sensibilidade) e \textit{mice} \cite{vanBuuren2011} (imputação). O nível de significância foi fixado em $\alpha = 0{,}05$, com correção de Bonferroni para múltiplas comparações dentro de cada dimensão.

\subsection{Outputs funcionais da meta-análise}
\label{sec:outputs_cap2}

A meta-análise foi projetada para produzir quatro entregas quantitativas funcionalmente integradas, que articulam o ranking de magnitudes de efeito ($\overline{lnRR}$) por dimensão ($V_1$--$V_6$), os índices de heterogeneidade ($I^2$, $\tau^2$) que sinalizam variabilidade contextual e condicionam transferibilidade, os coeficientes de meta-regressão dos moderadores contextuais (tipo de comunidade, região, tempo desde intervenção) e os intervalos de confiança de 95\% dos $lnRR$, permitindo hierarquizar a responsividade das dimensões a intervenções de salvaguarda e delimitar a precisão empírica das estimativas produzidas.

Para as dimensões $V_3$ e $V_5$ (estratégia exploratória), as estimativas são reportadas com cautela, sujeitas a atualização conforme novas evidências sejam incorporadas. Os outputs consolidados orientam investigações subsequentes, incluindo a construção futura de um índice composto de salvaguarda biocultural.

\subsection{Análise de lacunas de evidência}
\label{sec:gaps_cap2}

Para cada dimensão funcional ($V_1$--$V_6$), foi conduzida uma análise sistemática de lacunas de evidência, contabilizando o número de variáveis proxy listadas na Tabela~\ref{tab:proxies_cap2} para as quais nenhum estudo elegível reportou dados quantitativos suficientes para cálculo de $lnRR$. Essa contagem produz um vetor de ``evidência ausente'' por dimensão, que identifica as variáveis não cobertas pela literatura empírica existente e que, portanto, demandam elicitação em investigações complementares futuras. Adicionalmente, para cada variável proxy com evidência disponível, foi registrado se os dados provêm exclusivamente de contextos geográficos sem congruência com o Semiárido NE~II (moderador binário = 0), sinalizando variáveis cuja validade de transferência é condicionada e que requerem validação territorial complementar. A Tabela~\ref{tab:mapping_cap2} consolida o mapeamento entre os indicadores extraídos da meta-análise e as dimensões funcionais de vulnerabilidade, explicitando a correspondência entre outputs meta-analíticos e variáveis candidatas a compor instrumentos futuros de avaliação.

\begin{table}[htbp]
\caption{Mapeamento entre indicadores meta-analíticos e dimensões funcionais de vulnerabilidade biocultural.}
\label{tab:mapping_cap2}
\centering
\small
\begin{tabular}{p{0.8cm}p{2.0cm}p{3.5cm}p{2.5cm}p{3.0cm}}
\toprule
\textbf{Dim.} & \textbf{Denominação} & \textbf{Indicador meta-analítico ($lnRR$)} & \textbf{Variável funcional} & \textbf{Status de evidência} \\
\midrule
$V_1$ & Erosão Interg. & $\overline{lnRR}_{V_1}$, IC~95\% & Risco de descontinuidade de transmissão & Confirmatório ($k \geq 15$) \\
$V_2$ & Complexidade & $\overline{lnRR}_{V_2}$, IC~95\% & Densidade biocultural & Confirmatório ($k \geq 15$) \\
$V_3$ & Singularidade & $\overline{lnRR}_{V_3}$ (exploratório) & Exclusividade territorial & Exploratório ($k = 5$--$10$); requer validação futura \\
$V_4$ & Documentação & $\overline{lnRR}_{V_4}$, IC~95\% & Grau de registro formal & Confirmatório ($k = 8$--$12$); moderadores contextuais \\
$V_5$ & Vuln. Jurídica & $\overline{lnRR}_{V_5}$ (exploratório) & Exposição a apropriação indevida & Exploratório ($k < 8$); prioridade para elicitação futura \\
$V_6$ & Org. Social & $\overline{lnRR}_{V_6}$, IC~95\% & Vitalidade organizacional & Confirmatório ($k = 10$--$15$) \\
\bottomrule
\end{tabular}
\end{table}

A Figura~\ref{fig:methods-flow_cap2} sintetiza o encadeamento operacional completo, desde a definição do framework conceitual até a produção dos outputs funcionais da síntese meta-analítica.

\begin{figure}[H]
\centering
\includegraphics[width=0.82\textwidth]{CAPITULOS/ARTIGO_2/figuras/methods_flowchart.pdf}
\caption{Fluxograma metodológico integrando as etapas de escopo, triagem, classificação multi-tier, síntese estatística e produção de outputs funcionais.}
\label{fig:methods-flow_cap2}
\end{figure}


%% ============================================================
%% 3. RESULTADOS
%% ============================================================
\section{Resultados e Discussão}
\label{sec:results_cap2}

\subsection{Fluxo de seleção PRISMA}
\label{sec:prisma-flow_cap2}

A busca sistemática (2016--2026) recuperou 365 registros (193 Scopus via API, 172 Web of Science via exportação direta). Filtros de tipo documental e de janela temporal excluíram 33 registros não elegíveis (capítulos, conferências e editoriais) e 20 registros fora do período, reduzindo o corpus a 310 candidatos. A deduplicação cruzada por DOI e normalização de títulos eliminou 117 duplicatas interbases (sobreposição de 81\%), resultando em 193 registros únicos com contribuição líquida de 27 registros exclusivos da Web of Science.

A triagem semântica exigiu presença simultânea de descritores dos Blocos~1 (população) e 2 ou 3 (intervenção/desfechos), reduzindo o corpus de 193 para 42 registros. A correspondência fuzzy (\textit{stringdist}, método OSA, limiar $= 5$) via \textit{revtools} \cite{Westgate2019} não identificou duplicatas adicionais, porém a re-triagem manual de 139 registros sem resumo na exportação Scopus resgatou seis estudos elegíveis \cite{Aswani2020,Bussmann2018,Ghanimi2022,Malapane2024,RodriguezCruz2022,Rodriguez2018}, elevando o corpus de 42 para 48 registros e fixando a exclusão em 145.

A taxa de retenção de 24,9\% ($48/193$) reflete a especificidade do escopo, que exige convergência simultânea entre evidência populacional, mensuração de vulnerabilidade e intervenção de salvaguarda, combinação intrinsecamente rara na literatura indexada \cite{Agrawal2002,GomezBaggethun2013}. A exclusão dos 145 registros decorreu da convergência de ausência de desfechos quantitativos (perspectiva descritiva sem grupos de comparação), desalinhamento temático com o domínio de salvaguarda (blocos populacionais sem interseção com intervenção ou desfechos) e predominância de abordagens qualitativas sem parâmetros estatísticos extraíveis para cálculo de $lnRR$.

Os 48 registros retidos compreendem 46 artigos originais e duas revisões, publicados entre 2016 e 2026, com predomínio de registros Scopus ($n = 33$, 68,8\%) e contribuição complementar da Web of Science ($n = 15$, 31,2\%), conforme detalhado na Figura~\ref{fig:prisma-flow_cap2}.

\subsection{O paradoxo da vulnerabilidade reconhecida mas não mensurada}
\label{sec:thematic-profile_cap2}

O corpus de 48 registros concentra-se em \textit{traditional knowledge}, \textit{biocultural diversity}, \textit{agrobiodiversity} e \textit{ethnobotany}, confirmando a adequação da filtragem PICOS à interseção entre conhecimento ecológico local e conservação. Essa coesão temática mascara, contudo, uma disjunção epistemológica, pois virtualmente todos os registros reconhecem ameaças aos saberes tradicionais \cite{Mobarak2025,LaRosa2021,Suwardi2025,Essien2025}, porém apenas 7 dos 48 empregam ``vulnerabilidade'' como variável de investigação e somente \citeonline{Zeleke2025} operacionaliza o conceito com framework de 15 indicadores estruturados em exposição, sensibilidade e capacidade adaptativa.

A proporção de 31 estudos qualitativos para 17 quantitativos confirma que a vulnerabilidade biocultural permanece metodologicamente suboperacionalizada, reflexo de tradições etnográficas que tratam a quantificação com cautela por risco de descontextualização \cite{Agrawal2002}. A formulação condicional de \citeonline{Goncalves2022}, segundo a qual dados de uso de plantas alimentícias em comunidades quilombolas ``podem nos ajudar a compreender a vulnerabilidade à insegurança alimentar'', sintetiza o estado da arte, no qual a vulnerabilidade é percebida como horizonte interpretativo sem operacionalização paramétrica. A presente meta-análise constitui a ponte metodológica entre esse acervo descritivo e a quantificação necessária para subsidiar decisões de salvaguarda.

\subsection{Transmissão intergeracional como variável-nexo}
\label{sec:dimensional-coverage_cap2}

Dentre as seis dimensões de vulnerabilidade propostas ($V_1$--$V_6$), a transmissão intergeracional de saberes ($V_3$) se distingue na leitura do corpus não como uma dimensão equivalente às demais, mas como o mecanismo pelo qual todas as outras se reproduzem ou colapsam. Doze dos 48 estudos documentam explicitamente processos de transferência de conhecimento, e a diversidade dos contextos investigados revela que a transmissão não é um ato singular, mas um sistema de práticas distribuídas que assume formas radicalmente distintas dependendo do ecossistema, do gênero e da estrutura comunitária.

\citeonline{Cantor2025} demonstra que o conhecimento necessário para a pesca cooperativa com golfinhos na costa brasileira é transmitido por aprendizado situado, cuja ruptura implica não apenas perda de saber técnico, mas dissolução de uma relação interespecífica coevoluída. \citeonline{Sibanda2025} evidencia que a custódia feminina de sementes crioulas no Zimbabwe sustenta simultaneamente a soberania alimentar doméstica, a identidade cultural e a transmissão de conhecimento intergeracional, de modo que a vulnerabilidade seminal e a vulnerabilidade epistêmica são, nesse contexto, a mesma variável. \citeonline{Gerner2025} revela que a preservação de saberes sobre plantas medicinais na Europa depende de dinâmicas sociais comunitárias cujo enfraquecimento por migração desestrutura os mecanismos informais de transferência.

O que esses estudos demonstram em conjunto é que a transmissão intergeracional não é uma dimensão estanque, mas uma variável mediadora cuja falha precipita o colapso das demais dimensões. A co-ocorrência sistemática entre vetores de pressão exógena (migração, urbanização, integração de mercado) e erosão de saberes em pelo menos cinco estudos independentes \cite{Tran2025,Moller2023,Bastos2022,VazquezDelfin2022,Aniah2019} sustenta essa interpretação, na medida em que \citeonline{Tran2025} documenta como políticas de relocação no Vietnam desarticulam redes comunitárias de gestão hídrica tradicional, \citeonline{Moller2023} identifica a emigração como vetor primário de perda de saberes ecológicos em comunidades rurais sul-africanas, e \citeonline{Aniah2019} demonstra que a variabilidade climática no norte de Gana força a adoção de estratégias de subsistência alternativas que competem com a transmissão de práticas agrícolas tradicionais.

Essa centralidade funcional de $V_3$ tem implicação direta para a interpretação dos rankings de magnitude de efeito e para a eventual construção de instrumentos compostos de avaliação de vulnerabilidade. Se a transmissão intergeracional opera como nexo causal entre pressões exógenas e o decaimento de todas as dimensões, tratá-la com peso simétrico às demais subestimaria sua contribuição real para a vulnerabilidade agregada, o que deve ser considerado em estudos subsequentes de modelagem.

A dimensão de segurança alimentar ($V_4$) é documentada por seis estudos que convergem para uma interpretação funcional consistente, segundo a qual a diversidade de cultivos mantidos por sistemas tradicionais confere autonomia alimentar que protege comunidades contra flutuações de mercado e choques climáticos \cite{Sibanda2025,deSousa2024,Goncalves2022,Avila2021,Edo2024,RodriguezCruz2022}. \citeonline{deSousa2024} demonstra especificamente que a expansão da monocultura de soja sobre territórios quilombolas no Cerrado ameaça simultaneamente a base ecológica e o acervo biocultural dessas comunidades, produzindo um vetor duplo de vulnerabilidade que articula $V_4$ (segurança alimentar) com $V_2$ (complexidade biocultural).

A lacuna mais reveladora do corpus é a ausência de tratamento explícito da erosão cultural ($V_5$). Nenhum dos 48 estudos emprega esse termo como variável de investigação, embora processos equivalentes sejam descritos sob terminologias diversas, como ``erosão da diversidade biocultural'' \cite{Mobarak2025}, ``perda de conhecimento local e indígena'' \cite{Suwardi2025}, ``desaparecimento da história oral'' \cite{LaRosa2021}. Essa fragmentação terminológica indica que a erosão cultural é um fenômeno ubíquo mas desprovido de uma definição operacional consensual, o que motiva a classificação de $V_5$ como dimensão exploratória e a proposição de indicadores operacionalizáveis na presente síntese.

\subsection{Implicações para a avaliação de vulnerabilidade e direcionamento da extração}
\label{sec:vulnerability-implications_cap2}

A reflexão sobre o conteúdo dos registros retidos revela condicionantes que afetam diretamente a extração de dados e a interpretação dos resultados, demandando ajustes metodológicos antes da fase de leitura integral dos textos completos.

A composição metodológica do corpus (64,6\% de estudos qualitativos) reflete uma tensão epistemológica constitutiva do campo, na qual a quantificação de saberes tácitos e relacionais enfrenta objeções legítimas quanto a reificação \cite{Agrawal2002,GomezBaggethun2013}. A estratégia adotada nesta meta-análise responde parcialmente a essa tensão ao ancorar as estimativas em razões de resposta com propagação explícita de incerteza e ao manter $I^2$ como parâmetro de cautela interpretativa, preservando a gradatividade inerente ao fenômeno sem impor classificações dicotômicas.

Nessa perspectiva, o framework proposto por \citeonline{Suwardi2025}, que operacionaliza uma matriz de serviços ecossistêmicos vinculada a PCA e indicadores de transmissão de conhecimento, e a abordagem de \citeonline{Zeleke2025}, que emprega 15 indicadores quantitativos de vulnerabilidade, oferecem referências metodológicas diretamente transferíveis para investigações subsequentes que busquem construir instrumentos compostos de avaliação.

No plano geográfico, a concentração de estudos na América Latina (26,2\%) e África subsaariana (21,4\%) é favorável à transferibilidade dos parâmetros para comunidades quilombolas do Semiárido, na medida em que esses contextos compartilham sistemas agrícolas de subsistência em biomas tropicais com pressão fundiária crescente. Contudo, os quatro estudos brasileiros retidos \cite{Goncalves2022,Bastos2022,deSousa2024,Avila2021} foram conduzidos em biomas e territórios distintos do Semiárido Nordeste~II, condição controlada pelo moderador binário de congruência geográfica, cujos coeficientes de meta-regressão quantificam a transferibilidade dos parâmetros. A etapa complementar de validação in situ confrontará os parâmetros meta-analíticos com dados primários coletados em comunidades quilombolas do Vale do Catuxe.

A identificação de gênero como tema transversal em oito estudos (16,7\%) configura moderador funcional da vulnerabilidade biocultural. \citeonline{Sibanda2025} demonstra que a custódia feminina de sementes crioulas sustenta autonomia alimentar e transmissão de conhecimento simultaneamente, \citeonline{RomeroSilva2026} documenta o papel de mulheres na conservação de agrobiodiversidade em homegardens mesoamericanos, e \citeonline{MugiNgenga2016} identifica assimetrias de gênero na intensificação agrícola em comunidades quenianas. Esses achados sugerem que o moderador de gênero deve ser incorporado nas meta-regressões dimensionais para testar se comunidades com maior participação feminina na governança de saberes apresentam magnitudes de efeito distintas, particularmente em $V_3$ (transmissão) e $V_4$ (segurança alimentar), onde a evidência de mediação feminina é mais consistente.

\subsection{Limitações}
\label{sec:limitations_cap2}

A recuperação de registros abrangeu as bases Scopus e Web of Science Core Collection, não tendo sido incluídas bases regionais como SciELO Citation Index ou o Portal de Periódicos da CAPES, o que pode ter resultado na não captura de registros indexados exclusivamente nessas fontes. A sobreposição de 81\% entre os registros filtrados da Web of Science e o corpus Scopus sugere alta convergência de cobertura, porém a contribuição líquida de 27 registros exclusivos da Web of Science (14,0\% do corpus final) confirma a importância da busca multibase. A taxa de retenção de 24,9\% na triagem semântica permanece coerente com revisões sistemáticas que intersectam domínios disciplinares distintos (etnobiologia, ciências ambientais, patrimônio cultural) e indica que parcela expressiva da literatura indexada sobre saberes agrícolas tradicionais não reporta indicadores quantitativos de vulnerabilidade em formato extraível para meta-análise, reforçando a necessidade de elicitação primária de dados em investigações complementares futuras.

%% ============================================================
%% 5. CONCLUSÃO
%% ============================================================
\section{Conclusão}
\label{sec:conclusions_cap2}

A triagem inicial demonstrou a viabilidade de aplicar modelagem quantitativa à vulnerabilidade biocultural em saberes agrícolas tradicionais, revelando que 24,9\% da base filtrada (48 registros) apresenta aderência aos critérios rigorosos da arquitetura PICOS formulada. A eliminação justificada de 145 estudos confirmou o rigor do protocolo em separar menções puramente periféricas de desfechos agroecológicos e socioculturais mensuráveis. 

A arquitetura hierárquica estabelecida assegura que a extração subsequente fornecerá parâmetros matematicamente robustos (logaritmos da razão de resposta e seus respectivos intervalos de confiança de 95\%) livres de viés por agrupamento ou dependências intra-estudo. Os rankings de magnitude de efeito, os mapas de heterogeneidade e os coeficientes de moderadores produzidos constituem a base empírica que poderá subsidiar, em estudo subsequente, a construção de um índice composto de salvaguarda biocultural. A consolidação dos 48 registros elegíveis provê estimativas quantitativas que hierarquizam objetivamente as dimensões de vulnerabilidade e orienta a tomada de decisão em cenários de salvaguarda de comunidades quilombolas no Semiárido.


%% ============================================================
%% Declaração de contribuição dos autores (CRediT)
%% ============================================================
\section*{Declaração de contribuição dos autores (CRediT)}

\textbf{Catuxe Varjão de Santana Oliveira:} Conceituação, Metodologia, Investigação, Análise formal, Curadoria de dados, Escrita -- rascunho original. \textbf{Luiz Diego Vidal Santos:} Metodologia, Análise formal, Validação, Escrita -- revisão e edição.


%% ============================================================
%% Declaração de conflito de interesses
%% ============================================================
\section*{Declaração de conflito de interesses}

Os autores declaram não possuir conflitos de interesse financeiros ou relações pessoais conhecidas que possam ter influenciado o trabalho relatado neste artigo.


%% ============================================================
%% Disponibilidade de dados
%% ============================================================
\section*{Disponibilidade de dados}

O dataset de extração, os scripts de análise (R, versão 4.3+) e os outputs meta-analíticos parametrizados (CSV/JSON com metadados conforme princípios FAIR) estão disponíveis no repositório do projeto (\url{https://github.com/[repositorio-do-projeto]}). Dados brutos dos estudos primários estão sujeitos às licenças de publicação originais e podem ser obtidos mediante contato com os autores correspondentes dos estudos incluídos.


%% ============================================================
%% Agradecimentos
%% ============================================================
\section*{Agradecimentos}

O presente trabalho foi realizado com apoio da Coordenação de Aperfeiçoamento de Pessoal de Nível Superior -- Brasil (CAPES) -- Código de Financiamento 001. Os autores agradecem ao Programa de Pós-Graduação em Ciências da Propriedade Intelectual (PPGPI/UFS) pelo apoio institucional.


%% ============================================================
%% REFERÊNCIAS
%% ============================================================
